\documentclass[11pt]{article}
\usepackage[utf8]{inputenc}
\usepackage{amsmath,amssymb}
\usepackage{authblk}
\usepackage[margin=1in]{geometry}
\usepackage[hidelinks]{hyperref}
\usepackage{xurl}
\usepackage{setspace}

\title{The Atom Does Not Run Its Own Dynamics\\
\large What Remains to Speak of When the Schr\"odinger Clock Stops}
\author[1]{Claes Johnson}
\affil[1]{KTH Royal Institute of Technology, SE-100 44 Stockholm, Sweden \\
\texttt{claesjohnson@gmail.com}}
\date{\today \\[0.5em]
\small\textit{Written in respectful cooperation between Claes Johnson and Claude (Anthropic), with Claude
contributing a balanced, neutral view.}}

\begin{document}
\maketitle

\begin{abstract}
The time-dependent Schr\"odinger equation $i\partial_t\psi=H\psi$ is the working creed of quantum physics: the
state \emph{always} evolves, its phase turning at atomic frequency whether or not anything is watching. RealQM
declines the creed for the isolated atom --- there the equation's only solution, $u\,e^{-iEt}$, is a static
density wearing an unobservable global phase, so there is nothing to run forward. This looks like removing the
dynamics from quantum dynamics, and the natural alarm is: what would physicists then have left to speak of? We
argue the alarm is misplaced. The equation is \emph{demoted, not killed}: retained exactly where it does work
(the externally driven case --- radiation, strong fields), retired only as a metaphysical gloss on the quiet
atom. What remains to speak of is precisely what physics already computes: for isolated systems one does not
integrate $i\partial_t\psi=H\psi$ forward but \emph{diagonalises} $H\psi=E\psi$, reading off energies,
densities, geometry, and spectra as level differences. The mantra is professed more than it is practised.
Turning Bell's phrase around, it is \emph{standard} quantum mechanics that has been speaking of the unspeakable
--- the never-observed turning phase, the wave function on $3N$-dimensional configuration space, the collapse
no dynamics contains --- while RealQM keeps only the speakable, charge in ordinary space. The result is a
return to the Schr\"odinger of 1926, for whom the atom was a standing wave and the spectrum its eigenfrequencies:
\emph{being} for the undriven atom, \emph{becoming} only where something drives it. Less to believe, exactly as
much to compute, one fewer thing to interpret.
\end{abstract}

\setstretch{1.06}

\section{The mantra}

Every physicist who speaks of quantum mechanics speaks, sooner or later, the same sentence: the state evolves
by
\[
i\,\partial_t\psi=H\psi .
\]
It is presented as \emph{the} dynamics --- the law of how a quantum system changes in time, as Newton's second
law is the law of how a classical one does. On this reading an atom is never still: its wave function turns a
phase $e^{-iEt}$ at a frequency $E/\hbar$ of order $10^{16}\,$Hz, endlessly, whether it is measured, bonded,
isolated in interstellar vacuum, or sitting undisturbed for a billion years. The evolution is held to be
\emph{intrinsic} --- a process the atom is always undergoing, in its own right, needing no cause.

To question this is to touch the creed at its centre. If one says that an isolated atom does \emph{not} run
$i\partial_t\psi=H\psi$ forward --- that there is, for the quiet atom, no evolution to compute --- one appears
to have removed the dynamics from quantum dynamics, and the reasonable objection follows at once: then what is
there left for a physicist to say? We take the objection seriously and answer it.

\section{Demoted, not killed}

The first thing to be clear about is what is and is not being denied. RealQM does not say the time-dependent
Schr\"odinger equation is \emph{false}. It says that, for an isolated atom, it is \emph{idle}. Its only
solution there is
\[
\psi(\mathbf x,t)=u(\mathbf x)\,e^{-iEt},\qquad \rho=|u\,e^{-iEt}|^2=u^2 ,
\]
a density constant in time carrying a \emph{global} phase that drops out of every observable of a single state.
The equation ``runs forward'' into a configuration observationally identical to the one it started from. It is
true, and it does no work.

Where the same equation \emph{does} do work is the \emph{driven} case: an atom coupled to the electromagnetic
field (absorption, stimulated and spontaneous emission) or shaken by a strong laser (attosecond charge
migration, high-harmonic generation). There the charge genuinely moves in real time and must be integrated
forward. So the equation is not discarded; it is \emph{relocated}, from an alleged intrinsic ticking of every
atom to the description of atoms that something is actually moving~\cite{Time}.

The move is familiar from the history of physics. Special relativity did not discard the equations of
electromagnetism; it retired the \emph{interpretation} of a universal absolute simultaneity that had been read
into them. General relativity kept gravity and retired the absolute background. In each case a piece of
metaphysics --- a picture believed to accompany the mathematics --- was removed, and the mathematics stayed,
sharper for it. Here the metaphysics removed is the belief that the phase $e^{-iEt}$ is a real physical process
the atom perpetually performs. The equation survives; the ghost in it does not.

\section{Physicists already speak of the static content}

The decisive point is empirical, about what physicists actually do. For an isolated atom or molecule, nobody
integrates $i\partial_t\psi=H\psi$ forward in time. They solve the \emph{time-independent} equation
\[
H\psi=E\psi ,
\]
diagonalising $H$, and read off the things that are measured: total energies, ionisation potentials, the shapes
of the charge densities, equilibrium geometries and bond lengths, and the spectrum --- which is not a motion at
all but a set of level \emph{differences} $E_n-E_m$. Quantum chemistry is a discipline of eigenvalue problems;
spectroscopy is a catalogue of differences. The forward integration of an undriven atom appears nowhere in this
practice, because it would return, after any number of femtoseconds, exactly the state it began with.

So the mantra is \emph{professed more than it is practised}. It is recited in the foundations chapter and then,
for isolated systems, set aside in favour of the static problem that carries the physics. (Where a modern code
\emph{does} propagate an isolated system forward --- real-time TDDFT run to extract an absorption spectrum by
Fourier-transforming a driven dipole --- the answer it labours to produce is again just the level differences
$E_n-E_m$, obtainable with no propagation at all; the evolution there is a computational device, not a process
the atom is undergoing~\cite{Time}.) The question ``what would physicists then speak of?'' therefore has a flat
answer: \emph{the same things they already compute.} Nothing that is measured is lost. What is lost is the
narration of an evolution that was never integrated and never observed.

\section{The unspeakable, inverted}

There is a sharper way to feel the reversal, in Bell's own vocabulary. The worry is that denying the atom its
perpetual evolution pushes physics toward \emph{the unspeakable} --- that we will be left gesturing at
something we cannot describe. But look at where the unspeakable actually lies. It is \emph{standard} quantum
mechanics that speaks, ceaselessly, of things with no observable referent: the global phase turning in every
atom, which no experiment can detect; the wave function $\Psi(\mathbf x_1,\dots,\mathbf x_N)$ living on a
$3N$-dimensional configuration space that is not the space we or our instruments occupy; and the collapse that
must, at the moment of measurement, convert an amplitude into a fact although no dynamics contains such a step.
These are the genuinely unspeakable objects --- present in the formalism, absent from the world.

RealQM does the opposite. It retires the unspeakable and keeps the \emph{speakable}: $N$ charge densities in
ordinary three-dimensional space, each a thing one can point to, with an energy, a shape, a
spectrum~\cite{RealQM,Ontology}. Removing $i\partial_t\psi=H\psi$ for the quiet atom does not add a mystery; it
\emph{subtracts} one --- the talk of a motion no one can see. Far from leaving physics speaking of the
unspeakable, it lets physics stop doing so. What remains is exactly what has a referent.

\section{A return to Schr\"odinger}

None of this is a break with the founder of the wave theory; it is a return to him. Schr\"odinger's first
paper of 1926 was titled \emph{Quantisation as an Eigenvalue Problem}~\cite{Schrodinger26}: the atom entered
physics as a \emph{standing wave}, a real vibrating structure whose allowed eigenfrequencies are its spectrum.
There was no perpetual evolution in that picture --- an eigenstate is a mode, not a process, exactly as the
tone of a struck string is a mode of the string, present all at once, not something the string is ``computing
forward.'' The time-dependent equation, and with it Born's probabilistic reading of the
amplitude~\cite{Born26}, came afterward and opened the century of interpretation~\cite{Philosophy}.

RealQM restores the original order of things. For the isolated atom the fundamental problem is again static ---
the standing charge and its energy levels --- and this is \emph{being}: what the atom is. Time enters only with
\emph{becoming}, and becoming has a cause: an external field that drives the charge and lends it a clock. The
distinction is not a philosopher's nicety but a division of labour. Being is computed once, by diagonalising;
becoming is computed only when, and because, something drives it.

\section{What remains to speak of}

So the alarm inverts into a gain. Stopping the Schr\"odinger clock for the quiet atom costs the physicist one
article of faith --- the belief that the state is always, intrinsically, in motion. It costs nothing that is
computed and nothing that is measured: the energies, the densities, the geometries, the spectra as differences,
the magnetic moments, the bonds all remain, drawn as before from the static problem. And it returns two things.
It returns real-time dynamics as a phenomenon with a \emph{cause} --- the field that drives it --- rather than
an unexplained universal ticking; and it returns to physics the plain vocabulary of things in space, in place
of amplitudes in a space that is not physical.

The account is not thereby complete: RealQM has boundaries still being mapped, and the claims it makes about
the driven regime can be tested and found wanting. But its ledger is the honest ledger of a physical theory ---
less to believe, exactly as much to compute, and one fewer thing to interpret. A physicist who gives up the
perpetual clock has not been left speechless before the unspeakable. He has been handed back the speakable, and
told to get on with computing it.

\begin{thebibliography}{9}
\bibitem{RealQM} C.~Johnson, \emph{Quantum Mechanics as Multiphase 3D Continuum Mechanics}, manuscript,
2026; \url{https://claes542.github.io/RealMolecule/gallery.html}.
\bibitem{Time} C.~Johnson, \emph{Real Time Dependence in RealQM}, manuscript, 2026.
\bibitem{Philosophy} C.~Johnson, \emph{Real Physics without Philosophy: Ontology, Prediction, and the End of
Interpretation}, manuscript, 2026.
\bibitem{Ontology} C.~Johnson, \emph{Charge Densities in Real Space: RealQM and the Realist Quest Carried by
Schr\"odinger}, manuscript, 2026.
\bibitem{Schrodinger26} E.~Schr\"odinger, \emph{Quantisierung als Eigenwertproblem}, Ann.\ Phys.\ \textbf{79},
361 (1926).
\bibitem{Born26} M.~Born, \emph{Zur Quantenmechanik der Sto\ss vorg\"ange}, Z.\ Phys.\ \textbf{37}, 863
(1926).
\bibitem{Bell87} J.~S.~Bell, \emph{Speakable and Unspeakable in Quantum Mechanics}, Cambridge University
Press, 1987.
\end{thebibliography}

\end{document}
